\begin{abstract}
安全运营中心（SOC）收到的告警远多于分析师能调查的数量；无法把遥测数据发送给托管模型的组织，只能在自有硬件上用小型开源权重 LLM 来自动化分诊。让小型本地模型自己执行调查时，它们在流程上失败：不断探测却不收敛、从不给出结论，或把真实攻击判为良性。我们研究把调查流程移出模型之后会发生什么。\hesp{} 是一个控制器：它维护一个记录相互竞争的解释的账本，按“每单位成本的期望信息增益”挑选只读探针，只接受有当前证据支撑的结论，能自行结束调查，并在每次观测之前先把预测写入日志。我们用来自两个家族的五个开源权重模型（7B 到 72B）做了四项预注册研究，在我们自己搭建的受控分诊环境中共计 7{,}272 个经过审计的回合。把流程移出模型，使 Qwen2.5-7B 的已验证完成率从 0.125 提升到 1.000，使从不自行下结论的 Llama-3.1-8B 从 0 提升到 0.917；信息增益排序在每个会下结论的模型上带来 $+0.26$ 到 $+0.35$ 的提升。完全不含 LLM 的控制器本身就达到 0.917。在这个环境中，每个起因都会产生一个独特的离散观测，所以起作用的是流程而不是模型，模型可测的贡献只限于等到确认证据再停止。因此，我们的结果说明的是小型本地模型不应掌控调查流程，而不是控制器让它们成为更好的调查者；在观测模糊、必须从原始日志中读取的情形下，模型能否带来价值仍未检验。我们公开全部代码、协议和回合日志。
\end{abstract}
